\section{Literal canonical assignments and the repaired norm proof}
\label{sec:source}

\subsection{Typed terms and ordinal expression trees}

The ground type is $o$, and arrows associate to the right.  Put
\[
 \lv(o)=0,\qquad
 \lv(\sigma\to\tau)=\max\{\lv(\sigma)+1,\lv(\tau)\}.
\]
Terms of the original system T are variables, applications and
abstractions, together with
\[
 0:o,\quad S:o\to o,\quad
 D^\tau:o\to\tau\to\tau\to\tau,\quad
 R^\tau:o\to(o\to\tau\to\tau)\to\tau\to\tau.
\]
The defining reductions, in arbitrary well-typed contexts, are
\begin{align*}
 (\lambda X.A)B&\longrightarrow A[X:=B],\\
 D^\tau0AB&\longrightarrow A,&
 D^\tau(Sn)AB&\longrightarrow B,\\
 R^\tau0AB&\longrightarrow B,&
 R^\tau(Sn)AB&\longrightarrow A\,n\,(R^\tau nAB).
\end{align*}
Substitution is capture avoiding.  The assignment below reads the
literal typed term; it does not first normalize it or identify
extensionally equal terms.  The annotated recursor constructor used
temporarily in the reduction analysis of \cite{ww} is not a source
constructor here.

For each typed variable $X:\sigma$ introduce a vector
$\mathbf x=(x_0,\ldots,x_{\lv(\sigma)})$ of distinct ordinal
variables.  The ambient expression grammar has leaves $0,1,\omega$
and these component variables, and three binary constructors
\[
 \operatorname{Add}(a,b),\qquad
 \operatorname{ExpProd}(a,b),\qquad
 \operatorname{Psi}(a,b).
\]
Their values are respectively
\[
 a\oplus b,\qquad 2^a\otimes b,\qquad
 \psi_f(\omega\otimes a\oplus b).
\]
An exponential product or collapse is one primitive binary constructor,
not an arbitrary multiplication tree.  We often write $a+b$ for the
literal binary constructor $\operatorname{Add}(a,b)$; this shorthand
does not authorize reassociation.  Interpretation always uses natural
sum.  An expression is $\mathbf x$-free when no component of
$\mathbf x$ occurs in its tree.

A vector of level $m$ records exactly $m+1$ actual coordinate trees.
An out-of-range coordinate is absent and is interpreted as zero in a
comparison.  When vectors are added, two actual coordinates are joined
by one Add node, and a single actual coordinate is retained unchanged.
Thus addition does not insert actual Zero nodes to represent absent
coordinates.  All actual coordinates in the graded classes and in the
source assignment are zero-free and have positive values.

\subsection{The substitution domain and comparisons}

The closed classes $B_i$ are the least simultaneous classes satisfying
the following rules, which are those of \cite[Definition~2.6]{ww}.
\begin{enumerate}[label=(\roman*)]
\item $1\in B_i$ for every $i$, and $\omega\in B_i$ for $i\ge1$.
\item If $a,b\in B_i$, then $\operatorname{Add}(a,b)\in B_i$.
\item If $i\ge1$, $a\in B_{i+1}$ and $b\in B_i$, then
$\operatorname{ExpProd}(a,b)\in B_i$.
\item If $a\in B_1$, $b\in B_0$ and $N(a)\le F_2(b)$, then
$\operatorname{Psi}(a,b)\in B_0$.
\item Every member of $B_0$ is a member of $B_i$ for $i\ge1$.
\end{enumerate}
All expressions in $B_0$ have positive finite values.  A closed vector
$h$ belongs to $B$ when $h_i\in B_i$ at every actual coordinate.
It has \emph{bounded norm} when
\[
 N(h_i)\le N(h_0)=h_0\quad\text{for every }i.
\]

An admissible valuation replaces each occurring variable vector by one
whole vector in $B$ of the same level and of bounded norm.  If even one
component of a vector occurs in either expression being compared, all
components of that vector belong to the domain of the valuation.
Thus the components of a functional variable cannot be varied as
independent unconstrained ordinal parameters.

Scalar relations $a\mle b$, $a\mslt b$ and $a=b$ mean that their
values satisfy $\le$, $<$ and $=$ under every admissible valuation.
The same convention applies to norm inequalities.  For vectors use
the zero padding just specified and set
\[
 a\mle b\ \Longleftrightarrow\ (\forall i)\ a_i\mle b_i,
 \qquad
 a\mslt b\ \Longleftrightarrow\
 a_0\mslt b_0\ \text{ and }\ (\forall i>0)\ a_i\mle b_i.
\]
Equality is universal coordinatewise equality.  These are precisely the
comparison conventions of \cite[Definitions~2.7--2.8]{ww}.
For a ground variable the admissible values are exactly the positive
integers: every such integer is represented by a finite sum of ones,
and every $B_0$ expression has a positive integer value.

To describe the operators, define open classes $C_i$ by the same rules
as $B_i$, with the additional leaf $y_i\in C_i$ whenever
$i\le\lv(Y)$.  The collapse guard is now universal under admissible
valuations.  For a fixed $\mathbf x$, define $C_i^{\mathbf x}$ in
the same way, except that promotion from grade zero to a positive
grade is permitted only for a wholly $\mathbf x$-free expression.
Write $C$ and $C^{\mathbf x}$ for coordinatewise graded vectors.
Their closed instances are exactly $B$.  These classes do not include
a bounded-norm requirement; it is a separate property of a vector.

Induction on this grammar shows that an expression in
$C_i^{\mathbf x}$ contains no $x_j$ with $j<i$.
In particular, if $n=\lv(X)+1$, an expression of grade at least
$n$ is already $\mathbf x$-free.  This stopping property will pay
the norm budget in abstraction.

\subsection{A finite literal parsing policy}

Finite sums and repetitions in \cite[Definitions~2.10 and 2.18]{ww}
need an association convention to determine actual expression trees.
This choice matters when a later abstraction examines those trees.
For example the syntactic size defined below gives
\[
 s_{\mathbf x}(x_0+(1+1))=3,\qquad
 s_{\mathbf x}((x_0+1)+1)=5,
\]
although the two ordinal-valued expressions are equal.

\begin{definition}[Admissible literal policy]\label{def:literal-policy}
A policy $\rho$ is a computable choice of an ordered full binary
addition tree for every nonempty finite-sum macro site and arity.
It uses every operand once, in order and as an unchanged subtree.
Its choice depends on the site and arity, not on operand names,
values, internal trees or free-variable status.  A positive $k$-fold
repeat contains exactly $k$ unchanged copies; a one-fold repeat is
the operand itself.  Existing trees are never flattened or reassociated.
\end{definition}

The sites are positive numerals, positive repetitions, the active
three-operand sum partial, the active exponential partial, positive
coordinate aggregation in abstraction, its scalar repeat, and the
three-operand box tail when its finite contribution is nonzero.
In the exponential partial the first operand is the already formed
two-fold repeat of the head partial; it is not flattened into the
outer sum.  All expressly binary clauses remain binary.

One possible default uses left combs for the unspecified sum and repeat
sites, with box tail $a_0+(b_0+n)$ for $n>0$ and $a_0+b_0$ for
$n=0$.  The theorem allows any supplied policy in Definition~\ref{def:literal-policy}.
It neither requires nor asserts independence of the final assignment
from that policy.  A zero finite contribution in the box is absent;
this convention does not simplify any actual Zero tree.  The source
constant $0$ will be assigned the positive scalar $1$.

\subsection{Application and its domain estimate}

Let $a,b$ have levels $m>n$, and define $c=a\mathbin\square b$
downwards by
\begin{align*}
 c_i&=a_i&& (n<i\le m),\\
 c_i&=\operatorname{ExpProd}(c_{i+1},\operatorname{Add}(a_i,b_i))
                                      &&(1\le i\le n),\\
 c_0&=\begin{cases}
 \operatorname{Psi}(c_1,a_0+(b_0+n)),&n>0,\\
 \operatorname{Psi}(c_1,a_0+b_0),&n=0.
 \end{cases}
\end{align*}
The positive numeral $n$ and any optional association of the three
tail operands are expanded by $\rho$; their value is always
$a_0+b_0+n$.  Restriction $c\upharpoonright_r$ keeps the actual
coordinates $0,\ldots,r$.

\begin{lemma}[Box domain and bounded norm]\label{lem:source-box}
If $a,b\in C$ have bounded norm and $\lv(a)>\lv(b)=n$, then
$a\square b$ belongs to $C$ and has bounded norm.  The same statement
holds with $C^{\mathbf x}$ in place of $C$.
\end{lemma}
\begin{proof}
Fix an admissible valuation and put $u=a_0+b_0+n\ge n+2$.
All positive input-coordinate norms are at most $u$.
The product and exponential bounds in Lemma~\ref{lem:norm-arithmetic}
show that, at one downward step, a bound $x\ge u$ is replaced by
at most $u2^x$.  Since $u2^x\le2^{2^x}$ for $x\ge u\ge2$,
after at most $n$ steps every positive output norm is at most
$F_0^{2n}(u)$.  Directly from the iteration definitions,
\[
 F_0^{2n}(u)\le F_0^{2(u+1)}(u)\le F_1^2(u)\le F_2(u).
\]
For the middle inequality, expand both successive evaluations of
$F_1(v)=F_0^{v+1}(v)$; each contributes at least $u+1$
iterations beginning at a value at least $u$.

Consequently $N(c_1)\le F_2(u)$, which is the actual grammar guard
for $c_0=\operatorname{Psi}(c_1,u)$.  The positive coordinates
have their required grades by downward induction.  Put
$\alpha=\omega\otimes c_1\oplus u$.  It is nonfinite and
$N(\alpha)\ge u$, whence
\[
 c_0=\psi_f(\alpha)>f(N(\alpha))\ge G(u)\ge F_2(u).
\]
This bounds the norms of every positive output by its scalar output.
All estimates hold at every admissible valuation.  No promotion clause
is added in this construction, so the same reasoning proves the
$C^{\mathbf x}$ assertion.
\end{proof}

\subsection{The partials and the unchanged abstraction budget}

Fix $\mathbf x$ and put $n=\lv(X)+1$.  The syntactic size tests
whole-freeness first:
\[
 s_{\mathbf x}(h)=
 \begin{cases}
 1,&h\text{ is wholly }\mathbf x\text{-free, or }h=x_i,\\
 s_{\mathbf x}(a)+s_{\mathbf x}(b)+1,
    &h=\operatorname{Add}(a,b)\text{ or }\operatorname{Psi}(a,b),\\
 2s_{\mathbf x}(a)+s_{\mathbf x}(b)+1,
    &h=\operatorname{ExpProd}(a,b).
 \end{cases}
\]
The last two clauses apply only if the first does not.
For $h\in C_i^{\mathbf x}$, the vector $\partial_i^{\mathbf x}h$
has level $n$.  If $h$ is wholly $\mathbf x$-free, set
\[
 (\partial_i^{\mathbf x}h)_j=
 \begin{cases}h+1,&j=i\le n,\\1,&j\ne i\text{ or }i>n.
 \end{cases}
\]
For active expressions the clauses are
\begin{align*}
 \partial_i^{\mathbf x}x_i&=\mathbf1,\\
 \partial_i^{\mathbf x}\operatorname{Add}(a,b)
   &=\partial_i^{\mathbf x}a+\partial_i^{\mathbf x}b+\mathbf1,\\
 \partial_i^{\mathbf x}\operatorname{ExpProd}(a,b)
   &=2\partial_{i+1}^{\mathbf x}a+\partial_i^{\mathbf x}b+\mathbf1
       \qquad(i\ge1).
\end{align*}
All displayed vectors $\mathbf1$ have level $n$.  Multiples and
three-operand sums use the policy, including the nested two-fold head
repeat in the last clause.  For an active scalar collapse, put
\begin{align*}
 (\partial_0^{\mathbf x}\operatorname{Psi}(a,b))_0
   &=\operatorname{Psi}(a\{\mathbf x:=\mathbf1\},
                              (\partial_0^{\mathbf x}b)_0),\\
 (\partial_0^{\mathbf x}\operatorname{Psi}(a,b))_j
   &=(\partial_1^{\mathbf x}a)_j+(\partial_0^{\mathbf x}b)_j
                                      \qquad(1\le j\le n).
\end{align*}

For $h\in C^{\mathbf x}$ of level $m$, let
\[
 S_{\mathbf x}(h)=2^n\sum_{i=0}^m s_{\mathbf x}(h_i).
\]
This is a natural number computed from the fixed input trees.
The vector $\delta^{\mathbf x}h$ has level $\max\{n,m\}$ and
coordinates
\begin{equation}\label{eq:source-delta}
 (\delta^{\mathbf x}h)_j=
 \begin{cases}
 S_{\mathbf x}(h)\,(\partial_0^{\mathbf x}h_0)_0,&j=0,\\
 \displaystyle\sum_{i=0}^m(\partial_i^{\mathbf x}h_i)_j,&0<j\le n,\\
 h_j,&n<j\le m.
 \end{cases}
\end{equation}
The scalar multiplication is the literal positive repetition at its
designated policy site.  The natural number $S_{\mathbf x}(h)$ is
not recomputed after reassociating any input tree.

\begin{lemma}[Finite definition and substitution covariance]
\label{lem:source-finiteness}
For a supplied policy these definitions produce unique finite trees
effectively.  The partials and $\delta^{\mathbf x}$ commute with
substitutions for other variable vectors by $\mathbf x$-free vectors,
whenever the displayed operations have their graded domains.
\end{lemma}
\begin{proof}
Each partial is defined by structural recursion on its input tree.
For an explicit finite bound, let $|h|$ be its node count.  Give a leaf
weight one, Add and Psi weight $w(a)+w(b)+1$, and ExpProd weight
$2w(a)+w(b)+1$.  Induction gives
$s_{\mathbf x}(h)\le w(h)\le2^{|h|}-1$.
Every coordinate of a partial has at most $6\cdot3^{|h|}$ nodes:
the active Add, ExpProd, scalar Psi and positive Psi clauses have
bounds respectively
\[
 E(a)+E(b)+3,\quad2E(a)+E(b)+4,\quad
 |a|+E(b)+1,\quad E(a)+E(b)+1,
\]
where $E(a)=6\cdot3^{|a|}$, and the whole-free diagonal has
$|h|+2$ nodes.  Each fits the asserted recursive bound.
The finite sums and repeats in \eqref{eq:source-delta} therefore
have computable sizes.  This construction never has to decide a
universal comparison in order to construct a tree.

A substitution by $\mathbf x$-free vectors preserves whether a
subtree contains a component of $\mathbf x$.  Whole-free subtrees
still contribute size one, while the active constructor pattern is
unchanged.  Hence it preserves $s_{\mathbf x}$ and
$S_{\mathbf x}$.  Induction on each partial clause proves literal
commutation: the policy depends only on the unchanged macro site and
arity.  The finite sums in $\delta^{\mathbf x}$ then commute too.
\end{proof}

\begin{lemma}[Weak scalar estimate and the graded domain]
\label{lem:weak-scalar}
For $h\in C_0^{\mathbf x}$,
\[
 h\{\mathbf x:=\mathbf1\}\mle
             (\partial_0^{\mathbf x}h)_0.
\]
Every $\partial_i^{\mathbf x}h$, for $h\in C_i^{\mathbf x}$,
is an $\mathbf x$-free $C$-vector of level $n$.
For $h\in C^{\mathbf x}$, the vector $\delta^{\mathbf x}h$
is also valid and $\mathbf x$-free.
\end{lemma}
\begin{proof}
Induct simultaneously on the grammar and the partial recursion.
A whole-free scalar has diagonal $h+1$; the tracked scalar variable
has diagonal one; an active Add has the two scalar partials plus one.
For an active Psi the two heads after replacing $\mathbf x$ by ones
are identical, and the induction hypothesis compares their finite tails.
Lemma~\ref{lem:collapse-estimates} supplies the required weak
finite-tail monotonicity.

That new Psi has its grammar guard: the old guard and the weak
induction hypothesis give
\[
 N(a\{\mathbf x:=\mathbf1\})
 \mle F_2(b\{\mathbf x:=\mathbf1\})
 \mle F_2((\partial_0^{\mathbf x}b)_0).
\]
The head has the correct positive grade.  Whole-free diagonal outputs
are $h+1$ at their original grade; all off-diagonal and tracked-variable
outputs are one.  The other positive outputs only add correctly graded
partials, with the displayed fixed multiplicities.  Every output is
$\mathbf x$-free.  Sums and positive scalar repetitions preserve
these properties.  Above grade $n$, the retained $h_j$ is already
$\mathbf x$-free by the stopping property of the grammar.
\end{proof}

\begin{remark}[Two strict inequalities that are unnecessary]
The strict form printed in \cite[Lemma~2.19]{ww} fails at $h=x_0$:
both scalar sides are one.  Its appendix estimate (A.5) also has an
equality case: take $\lv(X)=0$, $h=(x_0,1)$, the coordinate
$i=1$, and its free subtree $1$.  The asserted strict estimate becomes
$1<1$.  The weak scalar lemma just proved is sufficient for the
collapse guard.  The next two results replace the strict free-subtree
estimate while preserving the multiplier in \eqref{eq:source-delta}.
These observations do not refute the separate beta estimate involving
both $\delta$ and $\square$ in \cite[Lemma~2.22]{ww}.
\end{remark}

\begin{lemma}[First free positive subtree]
\label{lem:first-free-subtree}
Let $h\in C^{\mathbf x}$ have bounded norm.  Fix an admissible
valuation for its other names and put
$B=h_0\{\mathbf x:=\mathbf1\}$ at that valuation.
In a partial recursion beginning at any actual $h_i$, every first
wholly $\mathbf x$-free subtree $a$ encountered at a positive
grade satisfies
\[
 N(a)\le2^{n-1}B.
\]
\end{lemma}
\begin{proof}
Measure norms after replacing $\mathbf x$ by ones.  The free
subtree is unchanged, and every initial coordinate then has norm at
most $B$.  A step into an Add child cannot increase norm.  A step
into an exponential tail cannot increase it either, by positivity and
the lower product bound.  A step into an exponential head costs at
most a factor two, since
\[
 N(a)\le2N(2^a)\le2N(2^a\otimes b).
\]
It also raises the grade by one.  At grade $n$ a subtree is already
wholly $\mathbf x$-free and recursion stops.  A path beginning
at positive grade thus makes at most $n-1$ such head steps.

A scalar path enters a positive grade only through a collapse head.
This step costs no factor two:
\[
 N(a)\le N(\omega\otimes a\oplus b)
       \le\psi_f(\omega\otimes a\oplus b).
\]
Its scalar tail is bounded by the same collapse value, so scalar
tail steps also cost no factor.  The head enters grade one, leaving
at most $n-1$ active exponential-head steps.  A scalar promoted to
positive grade is already wholly free, so it cannot begin another
active scalar descent.  Multiplication of the possible factors two
proves the bound.
\end{proof}

\begin{theorem}[Bounded norm with the original multiplier]
\label{thm:delta-bounded-norm}
If $h\in C^{\mathbf x}$ has bounded norm, then
$\delta^{\mathbf x}h$ has bounded norm, with exactly the multiplier
$S_{\mathbf x}(h)=2^n\sum_i s_{\mathbf x}(h_i)$ in
\eqref{eq:source-delta}, for every admissible policy.
\end{theorem}
\begin{proof}
Fix the other names as in Lemma~\ref{lem:first-free-subtree} and put
$H=(\partial_0^{\mathbf x}h_0)_0$ and $v=2^nH$.
The weak scalar estimate gives $1\le B\le H$.
Whenever a free positive-grade subtree can be copied into a positive
diagonal output, the preceding lemma gives
\[
 N(a)+1\le2^{n-1}B+1\le2^nH=v.
\]
A free \emph{scalar} subtree is not copied into a positive output;
those outputs are one.  This distinction supplies the extra unit in
the whole-free diagonal clause without a strict free-subtree bound.

Induct through the actual partial recursion, using those bounds at
each first free leaf.  For $0<j\le n$ the result is
\[
 N((\partial_i^{\mathbf x}h_i)_j)
           \le s_{\mathbf x}(h_i)v.
\]
Tracked-variable outputs have norm one.  An active Add contributes
the two child bounds and one, bounded by
$(s_{\mathbf x}(a)+s_{\mathbf x}(b)+1)v$.
An active ExpProd partial has two copies of its head output, one
tail output and one, bounded by
$(2s_{\mathbf x}(a)+s_{\mathbf x}(b)+1)v$.
A positive Psi output only adds its head and tail partials, so its
size allowance also suffices.  Norm additivity makes every allowed
association of these sums satisfy the same estimate.

Summing over $i$ gives
\[
 N((\delta^{\mathbf x}h)_j)
 \le\left(\sum_i s_{\mathbf x}(h_i)\right)2^nH
   =(\delta^{\mathbf x}h)_0\qquad(0<j\le n).
\]
For $j>n$ the retained $h_j$ is $\mathbf x$-free and its norm
is at most $B\le H\le(\delta^{\mathbf x}h)_0$.
Since the valuation was arbitrary, this proves bounded norm.
\end{proof}

\subsection{The canonical assignment and its arithmetic image}

The following is the zero-history instance of
\cite[Definition~3.1]{ww}, with literal expansions fixed by $\rho$:
\begin{align*}
 \can{X:\sigma}&=\mathbf x,\qquad
 \can0=(1),\qquad \can S=(1,1),\\
 \can{D^\tau}&=(1,\ldots,1)
                &&\text{of level }\lv(\tau)+1,\\
 \can{R^\tau}&=(2,1,\ldots,1,\omega)
                &&\text{of level }\lv(\tau)+2,\\
 \can{AB}&=(\can A\square\can B)\upharpoonright_{\lv(\tau)}
                &&(A:\sigma\to\tau, B:\sigma),\\
 \can{\lambda X.G}&=\delta^{\mathbf x}\can G+
                        \can G\{\mathbf x:=\mathbf1\}.
\end{align*}
The numeral $2$ is the literal $1+1$.  The vector of ones in a
substitution has the level of the variable being replaced.  When useful
we write $\can A_{f,\rho}$ to display the two fixed choices.

\begin{proposition}[Well-defined canonical source assignment]
\label{prop:canonical-domain}
Every literal original T term has an effectively constructed canonical
vector of the level of its type.  It belongs to $C$ and has bounded
norm.  A body to which $\delta^{\mathbf x}$ is applied belongs
to $C^{\mathbf x}$.  All these assertions allow arbitrary finite
internal and bound-variable types and arbitrary recursor instances.
\end{proposition}
\begin{proof}
The construction is structural recursion on the source term and the
finite recursions of Lemma~\ref{lem:source-finiteness}.  If
$A:\sigma\to\tau$, then
$\lv(A)\ge\lv(\sigma)+1>\lv(\sigma)$, which is exactly the
level premise for the box.  Restriction leaves level $\lv(\tau)$;
the abstraction has level $\max\{\lv(X)+1,\lv(G)\}$ as required.

For its domain, keep the following strengthened invariant: a positive
coordinate is generated from $1,\omega$, appropriate positive-index
component variables, Add and graded ExpProd, without any scalar
promotion or Psi.  The constants and variables have this property.
The positive box clauses preserve it.  In a scalar partial a whole-free
subtree contributes only ones to positive outputs, an active variable
contributes ones, and Add joins partials; an active Psi joins the
positive partials of its positive head and scalar tail.  In a partial
of a positive input, whole-free copies, variables, Add and ExpProd
also preserve the invariant.  The sums and retained coordinates in
abstraction do so as well.  This proves the invariant simultaneously
with source construction.  In particular every positive coordinate
has the grammar required for $C_i^{\mathbf x}$ for each tracked
binder, and its scalar collapse guards follow from
Lemmas~\ref{lem:source-box} and \ref{lem:weak-scalar}.

Variables have bounded norm by the valuation domain.  The constants
do too; for the recursor its top norm is $N(\omega)=2$, equal to
its scalar.  The box preserves bounded norm by
Lemma~\ref{lem:source-box}, as does restriction.  At an abstraction
the first summand has bounded norm by
Theorem~\ref{thm:delta-bounded-norm}; the second does because the
all-ones vector is admissible.  Coordinatewise natural addition
preserves the bound by norm additivity.  This completes the induction.
\end{proof}

\begin{theorem}[Fixed-head arithmetic image of actual sources]
\label{thm:generated-source}
Let $A$ be an original T term whose free variables all have ground
type.  For every positive coordinate, $\can A_i$ is a closed
effectively computable ordinal.  Its scalar is an effectively obtained
finite circuit in the ground component inputs, positive integer
constants, addition, fixed positive integer multiples, and functions
\[
 h_\lambda(u)=\psi_f(\lambda+u)
\]
with fixed nonzero limit heads $\lambda<\varepsilon_0$.
Repeated occurrences of one source variable give one shared input;
shared subcircuits may also be retained.  The theorem imposes no bound
on internal types, bound-variable types, result type, recursors, or
finite nesting, and is uniform in the supplied computable policy.
\end{theorem}
\begin{proof}
The positive-coordinate invariant in the preceding proof contains
only positive-index components of variables that remain free.
Every bound component is removed by the partials or by the simultaneous
all-ones substitution at its abstraction.  Ground free variables have
no positive-index components.  All positive output coordinates are
therefore closed, and their values are computable by the finite ordinal
operations of Lemma~\ref{lem:norm-arithmetic}.

The same induction shows that an actual scalar coordinate is formed
only from scalar component variables, positive constants, Add,
fixed repeats and Psi.  Every head of a scalar Psi is a positive
expression in the invariant just proved.  After the ground-free
restriction its value is a fixed closed ordinal $c>0$.
Set $\lambda=\omega\otimes c$.  Every monomial of this ordinal
has positive exponent, so $\lambda$ is a nonzero limit, and its
finite scalar tail $u$ gives
$\operatorname{Psi}(c,u)=\psi_f(\lambda+u)$.
The syntactic multiplier $S_{\mathbf x}(h)$ is a fixed integer
computed before valuations, so it adds no new variable or quantifier.
Translate this finite tree recursively, retaining names and shared
wires.  This is the required circuit.

Its head is the one produced by the literal source assignment.
The argument does not replace a generated head by an arbitrary ordinal
primitive or require a source term for every auxiliary head later used
by arithmetic elimination.
\end{proof}

The zero-history assignment is the object of comparison here.  The
normalization theorem in \cite{ww} permits history-dependent assignment
derivations at abstractions in order to follow reductions.  No claim
that this single fixed canonical map decreases at every reduction is
needed for the comparison theorem.
